\section{Выводы}

В этой лабораторной работе я реализовал словарь на собственном красно-чёрном дереве и довёл его до полного набора операций из условия: поиск, вставка, удаление, сохранение и загрузка. На практике самым важным оказалось не знание определения структуры, а аккуратная реализация изменяющих операций: нужно корректно менять связи между вершинами, обрабатывать крайние случаи и только после этого восстанавливать балансировку.

Самой содержательной частью работы оказалось удаление. Именно на этой операции стало понятно, насколько важны служебная вершина \texttt{nil\_}, корректные указатели на родителей и чёткое разделение между обычной BST-частью алгоритма и последующей балансировкой через \texttt{DeleteFixup}. Вставка устроена проще, но тоже требует аккуратности, чтобы после поворотов и перекрасок дерево оставалось в корректном состоянии.

Дополнительную практическую ценность дал бенчмарк. Он показал, что собственная реализация не только работает корректно, но и остаётся конкурентоспособной по сравнению с \texttt{std::map}: в смешанном и read-heavy сценариях результаты чаще были не хуже, а в сценарии с преобладанием изменений стандартная библиотека ожидаемо оказалась сильнее. Для меня это полезный итог, потому что он связывает теорию с практикой: важно не только реализовать структуру данных, но и понимать, как она ведёт себя на разных типах нагрузки.

Для меня основной результат этой лабораторной в том, что я не просто разобрал красно-чёрное дерево по теории, а реализовал его как рабочий класс на C++ и отработал детали, от которых реально зависит надёжность программы. В итоге получилась программа, которая корректно выполняет все операции словаря из задания и даёт понятное представление о том, где собственная реализация уже сопоставима со стандартной, а где ещё уступает ей по качеству исполнения.
\pagebreak
